%Chargement des paquets
\usepackage{amsmath}
\usepackage{amsfonts}
\usepackage{amssymb}
\usepackage{enumerate}
\usepackage{mathtools}
\usepackage{bbm}
\usepackage{xparse, etoolbox}
\usepackage{enumerate}
\usepackage{enumitem}
\usepackage{mathabx}
\usepackage{babel}[french]

%environment
\usepackage{amsthm}
\usepackage{keytheorems}
\usepackage{hyperref} 

%Interface théorème
\newkeytheorem{lem}[
	name = Lemme
]

\newkeytheorem{prop}[
	name = Proposition
]

\newkeytheorem{thm}[
	name = Théorème
]

\newkeytheorem{coro}[
	name = Corollaire
]

\renewcommand*{\proofname}{Démonstration}

\theoremstyle{definition}
\newtheorem{defn}{Définition}

\theoremstyle{definition}
\newtheorem*{nota}{Notation}

\theoremstyle{definition}
\newtheorem*{limi}{Liminaire}

\theoremstyle{remark}
\newtheorem{rmq}{Remarque}

\theoremstyle{plain}
\newtheorem*{fait}{Fait}


%Ensembles de nombres
\newcommand{\N} {\mathbb{N}}
\newcommand{\Ne}{\N^\ast}
\newcommand{\Z} {\mathbb{Z}}
\newcommand{\Q} {\mathbb{Q}}
\newcommand{\R} {\mathbb{R}}
\newcommand{\Rb}{\overline{\mathbb{R}}}
\newcommand{\Rp}{\R_+}
\newcommand{\Rm}{\R_-}
\newcommand{\K} {\mathbb{K}}
\newcommand{\Ket} {\mathbb{K^{\ast}}}
\newcommand{\Cx}{\mathbb{C}}

%Ensembles, tribus
\newcommand{\A}{\mathbb{A}}
\renewcommand{\P}{\mathcal{P}}
\newcommand{\T}{\mathcal{T}}
\newcommand{\F}{\mathcal{F}}
\newcommand{\C} {\mathcal{C}}
\newcommand{\ouv}{\mathcal{O}}
\newcommand{\B}{\mathcal{B}}
\newcommand{\M}{\mathcal{M}}
\newcommand{\etag}{\mathcal{E}}
\newcommand{\etagOp}{\etag^{+}(\Omega)}
\newcommand{\etagO}{\etag(\Omega)}
%%Lp avant quotientage
\newcommand{\Lic}{\mathcal{L}^1}
\newcommand{\Lpc}{\mathcal{L}^p}
\newcommand{\Linfc}{\mathcal{L}^{\infty}}
\newcommand{\Lifull}{\Lic(\Omega, \B, m)}
\newcommand{\Lpfull}{\Lpc(\Omega, \B, m)}
\newcommand{\Linffull}{\Linfc(\Omega, \B, m)}
%%Lp après quotientage
\newcommand{\Li}{L^1}
\newcommand{\Ld}{L^2}
\newcommand{\Lp}{L^p}
\newcommand{\Linf}{L^{\infty}}
\newcommand{\Listd}{\Li(\Omega, m)} 
\newcommand{\Ldstd}{\Ld(\Omega, m)} 
\newcommand{\Lpstd}{\Lp(\Omega, m)} 

%Quelques opérateurs
\DeclareMathOperator{\codim}{codim}
\DeclareMathOperator{\rg}{rg}
\DeclareMathOperator{\tr}{tr}
\DeclareMathOperator{\vect}{Vect}
\DeclareMathOperator{\ima}{Im}
\DeclareMathOperator{\id}{Id}
\DeclareMathOperator{\sign}{sign}
\DeclareMathOperator{\ext}{ext}
\newcommand{\ind}{\mathbbm{1}}
\newcommand*{\spr}[2]{\left(\,#1\,\middle|\,#2\,\right)}
\newcommand{\interior}[1]{%
  {\kern0pt#1}^{\mathrm{o}}%
}
\DeclareMathOperator{\card}{card}
\DeclareMathOperator{\ppcm}{ppcm}

%Commandes de pure flemme
\newcommand{\veps}{\varepsilon}
\newcommand{\intab}{\int_{a}^{b}}
\newcommand{\intR}{\int_{-\infty}^{\infty}}
\newcommand{\intinf}{\int_{- \infty}^{+ \infty}}
\newcommand{\equi}{\Leftrightarrow}
\newcommand{\Kev}{$\K$-espace vectoriel }
\newcommand{\Kevs}{$\K$-espaces vectoriels }
\newcommand{\Rev}{$\R$-espace vectoriel }
\newcommand{\Revs}{$\R$-espaces vectoriels }
\newcommand{\Cev}{$\Cx$-espace vectoriel }
\newcommand{\Cevs}{$\Cx$-espaces vectoriels }
\newcommand{\evn}{(E, \norm{.})}
\newcommand{\interf}[2]{[#1,#2]}
\newcommand{\limb}{\lim\limits_}
\newcommand{\lebn}{\lambda_n}
\newcommand{\pp}{\text{~presque partout}}
\newcommand{\derivp}[2]{\frac{\partial #1}{\partial #2}}
\newcommand{\famiu}{(u_i)_{i \in I}}
\newcommand{\sev}{\text{~s.e.v.}}
\newcommand{\Mnp}{M_{n,p}(\K)}
\newcommand{\Mn}{M_{n}(\K)}
\newcommand{\tq}{\text{~t.q.~}}
\newcommand{\mq}{~montrer~que~}
\newcommand{\et}{\quad \text{et} \quad}
\newcommand{\ou}{\text{~ou~}}
\newcommand{\Kx}{\K[X]}


%Commandes de gars chelou
\renewcommand{\subset}{\subseteq}

%Commande restriction, ty duckduckgo
\def\restr#1#2{\mathchoice
              {\setbox1\hbox{${\displaystyle #1}_{\scriptstyle #2}$}
              \restraux{#1}{#2}}
              {\setbox1\hbox{${\textstyle #1}_{\scriptstyle #2}$}
              \restraux{#1}{#2}}
              {\setbox1\hbox{${\scriptstyle #1}_{\scriptscriptstyle #2}$}
              \restraux{#1}{#2}}
              {\setbox1\hbox{${\scriptscriptstyle #1}_{\scriptscriptstyle #2}$}
              \restraux{#1}{#2}}}
\def\restraux#1#2{{#1\,\smash{\vrule height .8\ht1 depth .85\dp1}}_{\,#2}} 

%Quelques normes
\DeclarePairedDelimiter\abs{\lvert}{\rvert}
\DeclarePairedDelimiter\labs{\bigl\lvert}{\bigr\rvert}
\DeclarePairedDelimiter\norm{\lVert}{\rVert}
\DeclarePairedDelimiterXPP\normi[1]{}\lVert\rVert{_1}{\ifblank{#1}{\:\cdot\:}{#1}}
\DeclarePairedDelimiterXPP\normd[1]{}\lVert\rVert{_2}{\ifblank{#1}{\:\cdot\:}{#1}}
\DeclarePairedDelimiterXPP\normp[1]{}\lVert\rVert{_p}{\ifblank{#1}{\:\cdot\:}{#1}}
\DeclarePairedDelimiterXPP\normq[1]{}\lVert\rVert{_q}{\ifblank{#1}{\:\cdot\:}{#1}}
\DeclarePairedDelimiterXPP\norminf[1]{}\lVert\rVert{_\infty}{\ifblank{#1}{\:\cdot\:}{#1}}

%Bracket en angle
\DeclarePairedDelimiter\angl{\langle}{\rangle}

%Set, parenthèses
\newcommand{\set}[1]{\{ #1 \}}
\newcommand{\bigset}[1]{\bigl\{ #1 \bigr\}}
\newcommand{\Bigset}[1]{\Bigl\{ #1 \Bigr\}}
\newcommand{\bigpar}[1]{\bigl( #1 \bigr)}
\newcommand{\Bigpar}[1]{\Bigl( #1 \Bigr)}

%Opitmisation des opérateurs de comparaison
\DeclarePairedDelimiter\tio{o (}{)}
\DeclarePairedDelimiter\gro{O (}{)}


\newcommand{\Zpz}{(\Z/p\Z)^{\times}}
\newcommand{\Zpaz}{(\Z/p^{\alpha}\Z)^{\times}}
\newcommand{\D}{\mathcal{D}}

\pagestyle{fancy}

\fancyhead[C]{Cyclicité du groupe $\Zpaz$}
\fancyhead[R]{}

\begin{document}

\underline{Source :} Rombalid - Algèbre - page 292. \newline

\underline{Leçons :}
\begin{itemize}
\item 120 : Anneaux $\Z/n\Z. Applications.$
\item 121 : Nombres premiers. Applications.
\item 108 : Exemples de parties génératrices d’un groupe. Applications.
\item 104 : Groupes finis. Exemples et applications.
\end{itemize}

\begin{nota}
On note $p \geq 2$ un nombre premier et pour un entier $n \geq 1, \D_n$ l'ensemble des diviseurs de $n$. \\
Le symbole $\times$ en indice désigne l'ensemble des élements inversibles (pour la loi multiplicative) d'un anneau. \\
Pour tout $d \in \D_{p-1}$, on note $\psi(d)$ le nombre d'éléments d'ordre $d$ dans le groupe multiplicatif $\Zpz$. \\
La fonction indicatrice d'Euler est désigné par $\varphi$.
\end{nota}

\begin{lem}
On a, pour $\alpha \geq 2, \varphi(p^{\alpha}) = (p-1)p^{\alpha-1}$
\end{lem}

\begin{proof}
Notons $n$ un entier dans $\set{1, \dots, p^{\alpha}}$. Il y a en tout $p^{\alpha}$ tels entiers.
Parmi eux, seuls ceux dont l'écriture en base $p$ est de la forme $n = a_{\alpha-1}\dots a_1 a_0$ avec $a_0=0$ ne sont pas premiers
avec $p^{\alpha}$. Il existe $p^{\alpha-1}$ tels entiers (il s'agit de choisir une valeur dans $\set{0, \dots, p-1}$ pour chaque $a_i$,
exclure $n=0$ et inclure $n=p^{\alpha}$).
Finalement, on a bien $\varphi(p^{\alpha}) = p^{\alpha} - p^{\alpha-1} = (p-1)p^{\alpha-1}$.
\end{proof}

\begin{lem}
\label{lem:iso}
Tout groupe cyclique d'odre $n$ est isomorphe à $\left( \Z/n\Z, + \right)$.
\end{lem}

\begin{proof}
Soit $G = \angl{g}$ un tel groupe, considérons l'application $f : (\Z, +) \to G, k \mapsto g^k$,
on alors clairement $\ker{f}=n\Z$, et la restriction de $f$ à $\left( \Z/n\Z, + \right)$ induit un isomorphisme. 
\end{proof}

\begin{lem}
\label{lem:gen}
Soit $a \in \Z, \overline{a}$ est un générateur du groupe cyclique $\left( \Z/n\Z, + \right)$ si, et seulement si $a$ est premier avec $n$.
\end{lem}

\begin{proof}
$a$ est premier avec $n$ si, et seulement si, il existe $u, v \in \Z$ tels que $au+bv=1$.
Ce qui équivaut encore à $\overline{a}\overline{u}=\overline{1}$, donc $\overline{1}$ est dans le groupe additif généré par 
$\overline{a}$ ce qui équivaut à dire que ce groupe est $\Z/n\Z$ tout entier.
\end{proof}

\begin{rmq}
On a donc montré que $\left( \Z/n\Z, + \right)$ compte $\varphi(n)$ générateurs, ou encore que le groupe multiplicatif des inversibles
de $\Z/n\Z$ est d'ordre $\varphi(n)$.
\end{rmq}


\begin{lem}
Pour tout entier relatif $a$ premier avec $n$, on a $a^{\varphi(n)} \equiv 1 \pmod n$.
\end{lem}

\begin{proof}
Si $a$ est premier avec $n$, alors sa classe résiduelle est élement du groupe multiplicatif $(\Z/n\Z)^{\times}$, 
qui est d'ordre $\varphi(n)$, ce qui entraine directement le résultat (d'après le théorème de Lagrange). 
\end{proof}

\begin{lem}
On a $\psi(d) = \varphi(d)$ pour tout $d \in \D_{p-1}$.
\end{lem}

\begin{proof}
Si $\psi(d) >0$ alors il existe $x \in \Zpz$ d'ordre $d$ et le groupe $G = \set{\overline{1}, x, \dots, x^{d-1}}$ contient $d$ racines 
du polynôme $X^d - \overline{1}$ dans $\Z/p\Z$. Ce polynôme a au plus $d$ racines sur tout corps, donc $G$ est exactement l'ensemble
des racines de ce polynôme. Autrement dit, $G$ est l'ensemble des éléments d'ordre $d$ du groupe $\Zpz$.
Or, $G$ compte $\varphi(d)$ générateurs (c'est une conséquences des lemmes $\ref{lem:iso}$ et $\ref{lem:gen}$), 
ce qui montre bien que $\psi(d) = \varphi(d)$.
\end{proof}

\begin{coro}
Le groupe $\Zpz$ est cyclique.
\end{coro}

\begin{proof}
Le résultat est trivial pour $p=2$, sinon $p-1$ est pair et, selon le lemme précédent $\psi(p-1)=\varphi(p-1) > 0$, ce qui signifie 
qu'il existe au moins un élément d'ordre $p-1$ dans $\Zpz$ ce qui démontre que ce groupe est cyclique d'ordre $p-1$.
\end{proof}

\begin{thm}
Si $p$ est un nombre premier impair et $\alpha$ un entier supérieur ou égal à 2, alors le groupe multiplicatif $\Zpaz$ est cyclique.
\end{thm}

\begin{proof}
On démontre la suite de résultats suivants.
\begin{enumerate}[label=(\roman*)]

\item
$p$ divise tous les coefficients binomiaux $\binom{p}{k}$ pour $k \in \set{1, \dots, p-1}$, en effet :
\[ \binom{p}{k} = \dfrac{p \cdot (p-1) \cdots (p-k+1)}{k \cdot (k-1) \cdots 1} = \dfrac{a}{k!} \in \N , \]
On peut écrire $a = \binom{p^n}{k} k!$.
On remarque que $p$ et $k!$ sont premiers entre eux, or $p$ divise $a$, donc (théorème de Gauss), $p$ divise $\binom{p^n}{k}$.

\item
Il existe une suite d'entiers naturels non nuls $(a_k)_\N$ tous premiers avec $p$ et tels que :
\[ \forall k \in \N, \quad (1+p)^{p^k} = 1 + a_k p^{k+1} . \]
On construit la suite par récurrence sur $k \in \N$. Pour $k=0, a_0=1$ convient. \\
Pour $k=1$, on a :
\[ (1+p)^p = 1 + p^2 + \sum_{k=2}^p \binom{p}{k} p^k = 1 + p^2 ( 1 + \sum_{k=2}^p \binom{p}{k} p^{k-2} ) . \]
Or, chaque coefficient binomial est divisible par $p$ donc l'entier :
\[ a_1 = 1 + \sum_{k=2}^p \binom{p}{k} p^{k-2} \]
est non nul et premier avec $p$, il convient. \\
Supposons la suite construite jusqu'à un certains rang $k$ et on cherche à construire le terme $a_{k+1}$.
On a :
\begin{align*}
(1+p)^{p^{k+1}} & = \left ( (1 + p)^{p^k} \right )^p = \left ( 1 + a_k p^{k+1} \right )^p \\
 & = 1 + a_k p^{k+2} + \sum_{i=2}^p \binom{p}{i} a_k^i p^{(k+1)i} 
 & = 1 + p^{k+2} \left ( a_k + \sum_{i=2}^p \binom{p}{i} a_k^i p^{(k+1)i-2} \right ) .
\end{align*}
Comme pour le cas précédent, l'entier :
\[ a_{k+1} =  a_k + \sum_{i=2}^p \binom{p}{i} a_k^i p^{(k+1)i-2} \]
est non nul et premier avec $p$, donc convient. Finalement, la suite est constructible pour tout $n \in \N$, elle existe bien.

\item
La classe résiduelle de $1+p$ dans $\Zpaz$ est d'ordre $p^{\alpha-1}$. \\
Comme $1+p$ est premier avec $p^{\alpha}$, on a bien $\overline{1+p} \in \Zpaz$ et, selon le point (ii) :
\[ \left \{ \begin{array}{c c c}
(1+p)^{\alpha-1} & = 1 + a_{\alpha-1}p^{\alpha} & \equiv 1 \pmod p^{\alpha} \\
(1+p)^{\alpha-2} & = 1 + a_{\alpha-2}p^{\alpha-1} & \not \equiv 1 \pmod p^{\alpha}
\end{array} \right. \]
car $a_{\alpha-2}$ est premier avec $p$, donc $a_{\alpha-2}p^{\alpha-1}$ n'est pas divisible par $p^{\alpha}$.
On en déduit le résultat annoncé.

\item
Si $x = k+ p\Z$ est un générateur du groupe cyclique $\Zpz$, alors $y=k^{p^{\alpha-1}}+p^{\alpha}\Z$ est d'ordre $p-1$ dans
$\Zpaz$. \\
En effet, comme $x + p\Z$ génère $\Zpz$, sa classe résiduelle (tout comme celle de $k$) est d'ordre $p-1$ dans ce groupe. 
Par ailleurs, pour tout $\alpha \geq 2$, l'entier $p^{\alpha-1} - 1$ est divisible par $p-1$.
De plus $x^{p-1} \equiv k^{p-1} \equiv 1 \pmod p$, donc :
\[ k^{p^{\alpha-1}-1} \equiv k^{p-1} \equiv 1 \pmod p \]
et, par suite, $k^{p^{\alpha-1}} \equiv k \pmod p$. 
Autrement dit, la classe résiduelle de $j=k^{p^{\alpha-1}}$ est d'ordre $p-1$ dans $\Zpz$.
D'autre part :
\[ j^{p-1} = k^{(p-1)p^{\alpha-1}} = k ^{\varphi(p^{\alpha})} \equiv 1 \pmod p^{\alpha} , \]
Donc l'ordre de $y$ dans $\Zpaz$ divise $p-1$. Notons $r$ cet ordre, la relation :
\[ j^r \equiv 1 \pmod p^{\alpha} , \]
entraîne que $p^{\alpha}$, et donc $p$ divise $j^r-1$.
Autrement dit, $r$ est un multiple de $p-1$, il est donc égal à $p-1$.

\item
On peut enfin en déduire que $\Zpaz$ est cyclique. 
La classe résiduelle de $1+p$ dans $\Zpaz$ est d'ordre $p^{\alpha-1}$ (selon le point (iii)).
Il existe aussi un élément $y=k^{p^{\alpha-1}}+p^{\alpha}\Z$ dont la classe résiduelle 
dans $\Zpaz$ est d'ordre $p-1$ (selon le point (iv)).
Ainsi, l'odre de la classe résiduelle de $(1+p)y$ dans $\Zpaz$ est :
\[ \ppcm(p^{\alpha-1}, p-1) = (p-1) p^{\alpha-1} = \varphi(p^{\alpha}) , \]
On a donc montré que le groupe multiplicatif $\Zpaz$, possède un élément dont l'ordre est égal à son cardinal, 
le groupe est donc bien cyclique.
\end{enumerate}
\end{proof}




\end{document}